\sectionstart
\section{确定性现金计算与可解释规划}
\subsection{现金路径的精确定义}
金融计算使用整数分。金额入口拒绝布尔值、非有限数和超过分精度的输入，避免把浮点累计误差带入余额判断。设主账户已到账余额为 $B$，假设立即转出额为 $D$，每日演示开销为 $d$，第 $t$ 天的有符号事件金额之和为 $F_t$。则：
\begin{align}
C_0&=B-D,\\
C_t&=C_{t-1}-d+F_t,\qquad t=1,\ldots,30,\\
C_{\min}&=\min_{0\leq t\leq30}C_t,\qquad
G=\max\{0,R-C_{\min}\}.
\end{align}
其中 $R$ 为用户现金底线，$G$ 为预测底线缺口。路径包含今日和未来 30 天，共 31 个节点；最低值并列时取首次出现日。\code{free} 定义为 $C_{\min}-R$，可以为负值，不应直接称为可花余额。

事件包括既有账单、已登记的额外义务、有效订阅，以及尚未到账的预计工资。收到工资事件后，当前余额增加，同时移除预测中的这笔工资，防止重复计薪。延期改变预计到账日，移出观察窗口的收入不再计入 30 天路径。尚未确认的转账计划不是实际义务，不能因为模型提出建议就写入现金事实。

代码按每日开销和当日事件净额计算，只保存日终节点，不模拟同日先扣款后收入造成的日内缺口。第 0 天事件不在未来日循环中执行。当前演示账本与财务体检采用固定基准日，不能把输出称为已接入真实流水的滚动月报。
\source{backend/service.py: BankService._cash; backend/insights.py: build_insights}

\subsection{一组可以独立复算的合成算例}
设今日余额为 10,000 元，现金底线为 4,000 元，每日支出 100 元，第 5 天支付账单 3,000 元，第 10 天预计工资 6,000 元，第 12 天订阅扣费 100 元。以下均为说明算法而构造的输入，不是用户真实账本。

\begin{table}[htbp]\centering\small
\caption{同一账本在不同条件下的现金最低点}
\begin{tabularx}{\textwidth}{@{}YrrY@{}}\toprule
条件 & 最低现金（元） & 首次最低日 & 与底线比较 \\\midrule
原条件，不转出 & 6,100 & 第 9 天 & 高于底线 2,100 元 \\
原条件，立即转出 2,000 元 & 4,100 & 第 9 天 & 高于底线 100 元 \\
原条件，立即转出 2,500 元 & 3,600 & 第 9 天 & 低于底线 400 元 \\
工资延期 15 天，不转出 & 4,500 & 第 24 天 & 高于底线 500 元 \\
延期后，仍转出 2,000 元 & 2,500 & 第 24 天 & 低于底线 1,500 元 \\\bottomrule
\end{tabularx}
\end{table}
原条件下第 9 天余额为 $10000-9\times100-3000=6100$ 元；延期后第 24 天为 $10000-24\times100-3000-100=4500$ 元。相同转账额在两种条件下得到不同核验结果，体现状态变化对行动的约束。通过现金检查只说明这一组输入下的可行性，不产生用户授权，也不保证未来输入准确。

\subsection{财务体检的证据与规则}
体检复用服务端快照和现金路径，展示现金压力、小额支出、订阅复核项及其依据。单笔不超过 100 元且不属于内部调拨或还款的可见出账用于小额支出汇总；达到 2,000 元的出账或包含指令型文本的备注提示本人复核。低使用和重复订阅标记来自演示资料，不由模型假装推断真实使用频率。

现金低于用户底线时标记 \code{urgent}；没有现金缺口但存在复核信号时为 \code{attention}；没有上述信号时为 \code{steady}。这些是可解释规则状态，不是信用评分、欺诈判断或经过校准的风险概率。

取消单项订阅的比较，从窗口内对应续费日开始加回费用，不增加今日余额。由于最低点可能移动，多项取消对“最低余额改善”的影响不能把单项结果简单相加。体检同时列明数据、规则、未知项和决策边界；金融白话解释是只读知识入口，长辈模式只调整显示与关键确认信息的可读性。

\subsection{生活规划的有限离散搜索}
生活场景由用户声明金额、首选日、最晚日、重要性，以及是否允许替代或跳过。系统对保留原安排、缓冲优先、兼顾生活三种方案，比较基准、收入额外延期、突发支出、二者叠加四类压力条件。各情景没有估计概率；工资已到账时，延期情景可能与基准重合。

兼顾方案使用宽度为 128 的束搜索。候选日期来自首选日、最晚日以及允许窗口内的预计工资后日期，不逐日穷举；金额只保留用户原值、主动允许的替代金额，或在允许且非必要时跳过。必要场景不能自动取消。排序依次比较四情景下最坏支付缺口、最坏现金底线缺口、生活调整成本、节点缺口累计和日期指标；最终保留原安排与缓冲方案作为候选。

调整成本采用必要 100、重要 30、可调整 5 的演示权重，服务于可解释排序，不是估计效用或货币福利。有限束宽不保证全局最优。不动用资产只展示，不计入当前支付能力；缓冲敏感性固定同一方案与目标额，只改变初始现金假设。

采用生活安排会在事务中重算、复核双版本并登记未来义务，当前余额不扣款；采用不意味着所有压力情景均无风险。目标储蓄需另行生成转账草案，确认时走银行网关。生活分析、登记义务与实际付款具有不同状态，界面与回执必须分别表达。
\source{backend/life_planner.py: analyze, _choices, _rank, _policy; backend/life_service.py: commit_plan, propose_goal}
